\chapter{精选计算题详解}

本附录中的每题都重新写出条件和所求量，因此不依赖旧稿的“原专题”编号。计算从电子色散与场的归一化走到相干测量；每个近似都在使用处注明。

\section{自由电子与场的基本尺度}

\subsection*{相对论色散的两个导数}
电子正能量为 \(E(\mathbf p)=\sqrt{m^2c^4+c^2|\mathbf p|^2}\)。求群速度，并在一维纵向运动时求 \(E''(p_0)\)。

链式求导给
\[
\nabla_{\mathbf p}E=\frac{c^2\mathbf p}{E}.
\]
因此群速度就是粒子速度；在 \(p\ll mc\) 时，\(E=mc^2+p^2/(2m)+O(p^4)\)，于是 \(\mathbf v_g=\mathbf p/m+O(p^3)\)。再对纵向分量求导，
\[
E''(p)=\frac{c^2}{E}-\frac{c^4p^2}{E^3}
=\frac{m^2c^6}{E^3}=\frac1{\gamma^3m},
\qquad \gamma=\frac{E}{mc^2}.
\]
第二个导数控制窄动量波包的纵向色散；它不是 \(1/(\gamma m)\)，因为速度本身也随动量变化。

\subsection*{Dirac 正能旋量的流速}
设正能平面波旋量为
\[
u_s(\mathbf p)=N
\begin{pmatrix}\chi_s\\
\dfrac{c\boldsymbol\sigma\cdot\mathbf p}{E+mc^2}\chi_s
\end{pmatrix},
\qquad \chi_s^\dagger\chi_s=1.
\]
用 \(\boldsymbol\alpha=\begin{psmallmatrix}0&\boldsymbol\sigma\\
\boldsymbol\sigma&0\end{psmallmatrix}\) 计算 \(c\,u_s^\dagger\boldsymbol\alpha u_s/(u_s^\dagger u_s)\)。

记 \(B=c(\boldsymbol\sigma\cdot\mathbf p)/(E+mc^2)\)。Pauli 代数给
\[
u_s^\dagger\alpha_i u_s
=|N|^2\chi_s^\dagger(\sigma_iB+B\sigma_i)\chi_s
=|N|^2\frac{2cp_i}{E+mc^2}.
\]
同时 \((\boldsymbol\sigma\cdot\mathbf p)^2=p^2I\)，所以
\[
u_s^\dagger u_s
=|N|^2\left(1+\frac{c^2p^2}{(E+mc^2)^2}\right)
=|N|^2\frac{2E}{E+mc^2}.
\]
两式相除得到 \(c\,u_s^\dagger\alpha_i u_s/(u_s^\dagger u_s)=c^2p_i/E\)。归一化因子 \(N\) 消去；物理电荷流密度还要乘电荷 \(q=-e\)，但速度不依赖电荷符号。

\subsection*{相干态光子数的方差}
设单模相干态满足 \(a|\alpha\rangle=\alpha|\alpha\rangle\)，且 \(N=a^\dagger a\)。求 Fock 系数、平均光子数及方差。

若 \(|\alpha\rangle=\sum_{n\ge0}c_n|n\rangle\)，逐项比较 \(a|\alpha\rangle\) 与 \(\alpha|\alpha\rangle\) 得 \(c_{n+1}\sqrt{n+1}=\alpha c_n\)。归一化后
\[
|\alpha\rangle=e^{-|\alpha|^2/2}
\sum_{n=0}^\infty\frac{\alpha^n}{\sqrt{n!}}|n\rangle.
\]
令 \(\bar n=|\alpha|^2\)。由 \(N(N-1)=a^{\dagger2}a^2\) 及本征态性质，
\[
\langle N\rangle=\bar n,\qquad
\langle N(N-1)\rangle=\bar n^2.
\]
于是 \(\langle N^2\rangle=\bar n^2+\bar n\)、\(\operatorname{Var}N=\bar n\)。平均能量相同的数态方差为零，热态方差为 \(\bar n(\bar n+1)\)；只给平均光子数不足以确定交换统计。

\subsection*{Green 张量的量纲与点偶极源}
频域约定为 \(e^{-\ii\omega t}\)，并定义电场的推迟 Green 张量，使
\[
\mathbf E(\mathbf r,\omega)=\ii\omega\mu_0
\int_{\mathbb R^3}\mathbf G(\mathbf r,\mathbf r';\omega)
\mathbf j(\mathbf r',\omega)\dd^3r'.
\]
求 \(\mathbf G\) 的 SI 量纲；再将位于 \(\mathbf r_0\) 的点偶极电流 \(\mathbf j=-\ii\omega\mathbf d\,\delta(\mathbf r'-\mathbf r_0)\) 代入。

\(\mathbf j\,\dd^3r'\) 的单位为 \(\si{\ampere\meter}\)，\(\omega\mu_0\) 的单位为 \(\si{\ohm\per\meter}\)，两者相乘是 \(\si{\volt}\)。为了得到电场的 \(\si{\volt\per\meter}\)，\(\mathbf G\) 必须具有 \(\si{\per\meter}\) 的量纲。代入点源并使用 delta 函数的抽样性质，得到
\[
\mathbf E(\mathbf r,\omega)=\mu_0\omega^2
\mathbf G(\mathbf r,\mathbf r_0;\omega)\mathbf d.
\]
若求偶极对自身场的响应，需要适当处理 \(\mathbf r=\mathbf r_0\) 的奇异自场；移动电子的电流分布则会在两点 \((\mathbf r,\mathbf r')\) 上采样 Green 核，不能简单替成单个点偶极矩阵元。

\section{有限作用时间与边带}

\subsection*{两种时间窗的失谐响应}
给定失谐 \(\Delta\)，分别计算矩形窗 \(w_R(t)=\mathbf1_{[-\tau/2,\tau/2]}(t)\) 与高斯窗 \(w_G(t)=e^{-t^2/(2\tau^2)}\) 的振幅 \(A(\Delta)=\int_{\mathbb R}w(t)e^{\ii\Delta t}\dd t\)。

直接积分与配方分别给
\[
A_R(\Delta)=\tau\frac{\sin(\Delta\tau/2)}{\Delta\tau/2},
\qquad
A_G(\Delta)=\sqrt{2\pi}\tau e^{-\Delta^2\tau^2/2}.
\]
矩形窗在 \(\Delta\tau=2\pi m\)、\(m\in\mathbb Z\setminus\{0\}\) 时严格为零；高斯振幅对任意有限实 \(\Delta\) 都为正。两窗的参数 \(\tau\) 定义不同，比较线宽前必须说明是固定峰值、面积、均方时宽还是其他尺度，不能仅说“长度相同”。

\subsection*{反冲使能量梯不再等间距}
一维电子初动量为 \(p_0\)，相互作用后第 \(\ell\) 条动量边带为 \(p_\ell=p_0+\ell\hbar k_z\)。假设 \(|\ell\hbar k_z|\ll |p_0|\) 且所用边带范围内三阶 Taylor 余项可忽略。求相邻两条能量差的差值至 \(k_z^2\) 阶。

在 \(p_0\) 处展开，
\[
E_\ell=E_0+\ell\hbar k_zv_0
+\frac{\ell^2\hbar^2k_z^2}{2\gamma_0^3m}
+O((\ell\hbar k_z)^3),
\qquad v_0=E'(p_0).
\]
故
\[
E_{\ell+1}-E_\ell
=\hbar k_zv_0
+\frac{(2\ell+1)\hbar^2k_z^2}{2\gamma_0^3m}
+O((\ell\hbar k_z)^3),
\]
而连续两项相减得到 \(\hbar^2k_z^2/(\gamma_0^3m)\) 至所示阶数。这是能量梯非均匀性的尺度；若它乘作用时间再除以 \(\hbar\) 已不可忽略，就不能把所有边带当作同一共振频率。

\subsection*{弱场 PINEM 的概率守恒}
在经典场、低反冲且理想相干模型中，边带概率为 \(P_\ell=J_\ell^2(2|\beta|)\)，其中 \(\beta\) 是无量纲复耦合。令 \(b=|\beta|\ll1\)，求 \(P_0,P_{\pm1}\) 至 \(b^2\) 阶。

Bessel 函数在原点的级数给 \(J_0(2b)=1-b^2+O(b^4)\) 与 \(J_{\pm1}(2b)=\pm b+O(b^3)\)。因而
\[
P_0=1-2b^2+O(b^4),\qquad
P_{+1}=P_{-1}=b^2+O(b^4).
\]
三项相加为 \(1+O(b^4)\)；更高边带的概率从 \(b^4\) 阶才开始贡献。中心边带的减少同时供给两个一阶边带，不能分别给它们指定互不相关的概率。

\subsection*{横向束斑对局域耦合的平均}
电子在横向平面的归一化概率密度为
\(w(R)=(2\pi\sigma_R^2)^{-1}e^{-R^2/(2\sigma_R^2)}\)，局域耦合为
\(\beta(R)=\beta_0e^{-R^2/(2L_E^2)}\)。在弱耦合 \(P_1(R)=|\beta(R)|^2+O(|\beta|^4)\) 下求平均概率的主阶。

使用极坐标面积元 \(2\pi R\dd R\)，有
\[
\overline P_1
=|\beta_0|^2\int_0^\infty
\frac{R\dd R}{\sigma_R^2}
e^{-R^2/(2\sigma_R^2)-R^2/L_E^2}
=\frac{|\beta_0|^2}{1+2\sigma_R^2/L_E^2}.
\]
当 \(\sigma_R/L_E\to0\) 时恢复轴上单电子结果；束斑远宽于近场时，多数电子采样的耦合很弱，平均信号被压低。这里平均的是概率而非振幅，因为未分辨电子的横向入射位置。

\subsection*{损失峰与发光峰为何可以不重合}
设电子经过一个线性吸收结构附近，某模式从电子获得能量的速率为 \(\Gamma_{\rm exc}(\omega)\)。令进入远场光子的分支比为 \(\eta_{\rm rad}(\omega)\)，探测器收集其中的比例为 \(\eta_{\rm coll}(\omega)\)。写出该模式对收集发光的贡献，并解释为何能损有峰而发光可以没有峰。

在可以把该模式与其他模式区分、且每次激发的分支统计适用时，收集发光率为
\[
\Gamma_{\rm coll}(\omega)
=\Gamma_{\rm exc}(\omega)\eta_{\rm rad}(\omega)
\eta_{\rm coll}(\omega).
\]
电子能损统计电子交出的能量，不要求这份能量最终成为光子；发光则多了辐射分支和孔径投影。若一个模式把总能量的 \(80\%\) 辐射到远场，而探测器收集其中 \(25\%\)，总收集比例是 \(0.8\times0.25=0.2\)。改变数值孔径可以改变 \(\eta_{\rm coll}\)，不能反向改变结构内部已发生的吸收。对真正的多模干涉或非线性过程，简单分支比乘法需让位于对完整输出通道的计算。

\section{联合态、传播噪声与测量}

\subsection*{偏迹与条件态是两种不同操作}
设归一化电子--光联合态为
\[
|\Psi\rangle=\frac{|0\rangle_e|0\rangle_\gamma+
|1\rangle_e|1\rangle_\gamma}{\sqrt2}.
\]
求未观测光场时的电子态；再把光投影到 \(|\pm\rangle_\gamma=(|0\rangle_\gamma\pm|1\rangle_\gamma)/\sqrt2\)，求两个条件电子态及其概率。

对光作偏迹时，交叉项因 \(\langle0|1\rangle_\gamma=0\) 消失，故
\[
\rho_e=\tfrac12(|0\rangle\langle0|+|1\rangle\langle1|).
\]
投影到 \(|\pm\rangle_\gamma\) 后，未归一化电子矢量为
\((|0\rangle_e\pm|1\rangle_e)/2\)，范数平方均为 \(1/2\)，归一化后分别为 \(|\pm\rangle_e\)。若不保留光测量结果，则
\[
\tfrac12|+\rangle\langle+|+\tfrac12|-\rangle\langle-|
=\rho_e.
\]
后选择可以在某个子样本中恢复相干，但不能凭空改变未选择样本的态。

\subsection*{指数相关相位噪声的两个时间极限}
设实平稳零均值 高斯 噪声 \(\xi(t)\) 满足 \(C(u)=\mathbb E[\xi(t+u)\xi(t)]=\sigma^2e^{-|u|/\tau_c}\)。令 \(\Phi(t)=\int_0^t\xi(s)\dd s\)，求 \(\operatorname{Var}\Phi(t)\) 并比较 \(t\ll\tau_c\) 与 \(t\gg\tau_c\)。

将方形积分域按 \(s>s'\) 与 \(s<s'\) 分开，再以 \(u=|s-s'|\) 换元，
\[
\operatorname{Var}\Phi(t)=
\int_0^t\!\int_0^t C(s-s')\dd s\dd s'
=2\int_0^t(t-u)C(u)\dd u
=2\sigma^2\bigl[t\tau_c-\tau_c^2(1-e^{-t/\tau_c})\bigr].
\]
短时展开为 \(\sigma^2t^2+O(t^3/\tau_c)\)，长时为 \(2\sigma^2\tau_ct+O(1)\)。若相干因子为 \(\mathbb E e^{\ii\Phi}\)，高斯 特征函数给它等于 \(e^{-\operatorname{Var}\Phi/2}\)：起初是 高斯 型，只有长时才是近似指数型。

\subsection*{有限长度光栅的相位匹配宽度}
设一维光栅沿 \(z\) 方向周期为 \(a\)，选取非零整数级次 \(m\) 的倒格矢 \(G_m=2\pi m/a\)。电子速度为 \(v>0\)，目标辐射频率为 \(\omega>0\)。在所选光学通道纵向波数 \(k_z=0\) 的简化情形，相位失谐为 \(\Delta k=\omega/v-G_m\)。求精确匹配的周期，并说明长为 \(L\) 的均匀相互作用区如何展宽条件。

精确匹配要求 \(\Delta k=0\)，即 \(a=2\pi mv/\omega\)；由于 \(a>0\)，在当前符号约定下必须选 \(m>0\)。有限区间的振幅因子为
\[
\int_{-L/2}^{L/2}e^{\ii\Delta k z}\dd z
=L\,\frac{\sin(\Delta kL/2)}{\Delta kL/2}.
\]
因此匹配条件不再是严格的 delta 线；中心峰的第一零点位于 \(|\Delta k|=2\pi/L\)，典型失谐宽度为 \(O(L^{-1})\)。实际光学通道若有非零 \(k_z\)，应退回完整条件 \(\omega/v=k_z+G_m\)，而不能沿用这里的特殊式。

\subsection*{三个参考相位怎样识别复耦合}
在弱耦合测量中，设信号为 \(u=Ae^{\ii\phi}\)，已知参考幅值 \(B>0\)，在设置 \(\theta_s\) 下可观测量为 \(p_s=|u+Be^{\ii\theta_s}|^2\)。取 \(\theta_s=0,\pi/2,\pi\)，判断能否从三个理想无噪数据确定 \(u\)。

展开平方得
\[
p_\theta=A^2+B^2+2AB\cos(\phi-\theta).
\]
令 \(x=A\cos\phi\)、\(y=A\sin\phi\)，则
\[
p_0-p_\pi=4Bx,\qquad
p_{\pi/2}-\frac{p_0+p_\pi}{2}=2By.
\]
由于 \(B>0\)，两式分别给出 \(x\) 与 \(y\)，也就给出复数 \(u=x+\ii y\)。当 \(A=0\) 时 \(u=0\) 仍可识别，但极坐标参数 \(\phi\) 没有定义；这时说 \((A,\phi)\) 的 Jacobian 失秩，是参数化在原点奇异，并不表示直角坐标 \((x,y)\) 不可识别。若仪器能量分辨率使这些理想 \(p_s\) 无法区分，还必须把响应函数纳入前向模型，代数可逆性本身并不保证实验可识别。
